\documentclass{article}
\usepackage[letterpaper,portrait,margin=1in]{geometry}
\usepackage{amsmath}
\usepackage{amsthm}
\usepackage{fancyhdr}
\usepackage{lastpage}

\pagestyle{fancy}
\fancyhf{} %Clear the header and footer (don't use the default plain page style).
\renewcommand{\footrulewidth}{0.4pt}
\lhead{CLASS -- \textbf{Assignment Name}} %Print text on the left side of the header.
\rhead{Author Name} %Print text on the right side of the header.
\rfoot{Page \thepage \ of \pageref{LastPage}} %Print the page on the right side of the footer.
%Optional: Can print text in the center of the header with \chead{}.
%\chead{Center test}

\title{\vspace*{-1cm} CLASS NAME \\* \textbf{Assignment Name}}
\author{Author Name \\* \\* \textit{Collaborators: Names}}
\date{Due date, or some other relevant date}

%Define commands
\newtheorem*{claim}{Claim}

\theoremstyle{definition}
\newtheorem*{example}{Example}

\theoremstyle{plain}
\newtheorem*{lemma}{Lemma}

\newtheorem*{remark}{Remark}

\begin{document}

\maketitle
\thispagestyle{fancy}

%Reminder: Remove the * at the end for automatic numbering.
\section*{Problem 13.3 of A\&A}
\subsection*{Part a)}
\begin{claim}
\(\sqrt{2}\) is irrational.
\end{claim}
\begin{proof}
Suppose that \(\sqrt{2}\) were rational. There there exist some non-zero integers \(p\) and \(q\) such that \[\sqrt{2} = \frac{p}{q}.\] Without loss of generality, we can assume that this fraction is simplified so that \(p\) and \(q\) have no common divisors. Then,
\begin{align*}
    \sqrt{2}q &= p \\
    2q^{2} &= p^{2}.
\end{align*}
This implies (by the fundamental theorem of arithmetic) that 2 divides \(p^{2}\) and, hence, 2 divides \(p.\) Therefore, \(p = 2r\) for some (non-zero) integer \(r.\) Therefore,
\begin{align*}
    2q^{2} &= p^{2} \\
    2q^{2} &= (2r)^{2} \\
    2q^{2} &= 4r^{2} \\
    q^{2} &= 2r^{2}.
\end{align*}
From this, we conclude that 2 divides \(q^{2}\) and, hence, 2 divides \(q.\) This is a contradiction as we assumed \(p\) and \(q\) were written so that they share no common divisors (other than the trivial divisor 1, of course). Therefore, \(\sqrt{2}\) is not a rational number.
\end{proof}

\subsection*{Part b)}
To prove the claim of the problem, note that...
\end{document}